\chapter{Introdução}
\label{cap:introducao}

As doen\c cas cardiovasculares representam a principal causa de mortalidade no Brasil e no mundo. Segundo a Diretriz Brasileira de Insufici\^encia Card\'iaca, as doen\c cas do sistema circulat\'orio respondem por aproximadamente 30\% dos \'obitos no pa\'is, com tend\^encia de crescimento associada ao envelhecimento populacional, \`a urbaniza\c c\~ao acelerada e \`a preval\^encia cada vez maior de fatores de risco cardiovascular, como hipertens\~ao arterial, diabetes e dislipidemia \cite{rohde2018}. A insufici\^encia card\'iaca, em particular, afeta mais de 2 milh\~oes de brasileiros, com estimativas de que metade dos diagnosticados n\~ao sobreviva aos cinco anos seguintes ao diagn\'ostico \cite{rohde2018}.

Diante desse quadro epidemiol\'ogico, o diagn\'ostico precoce emerge como interven\c c\~ao de maior impacto na redu\c c\~ao da mortalidade e complica\c c\~oes cardiovasculares. Identificar as doen\c cas cardiovasculares nas fases iniciais, antes da instala\c c\~ao de complica\c c\~oes graves como o infarto do mioc\'ardio e o acidente vascular encef\'alico, amplia significativamente as op\c c\~oes terap\^euticas dispon\'iveis e melhora o progn\'ostico do paciente \cite{porto2019}. No entanto, o diagn\'ostico preciso de doen\c cas card\'iacas frequentemente depende de exames especializados, como ecocardiografia, cineangiocoronariografia e testes de esfor\c co, cuja disponibilidade \'e desigualmente distribu\'ida no territ\'orio nacional, especialmente em regi\~oes com menor infraestrutura assistencial \cite{brandao2025}.

\'E nesse contexto que modelos preditivos baseados em dados cl\'inicos acess\'iveis, como idade, press\~ao arterial, colesterol, eletrocardiograma de repouso e hist\'orico de sintomas, passam a ter relev\^ancia pr\'atica crescente. Um algoritmo capaz de estimar, a partir dessas vari\'aveis de baixo custo, a probabilidade de doen\c ca card\'iaca em um paciente poderia apoiar decis\~oes cl\'inicas em ambientes de aten\c c\~ao prim\'aria, orientar a prioriza\c c\~ao de exames avan\c cados e contribuir para a triagem em larga escala de popula\c c\~oes de risco \cite{mesquita2020}.

\section{Contextualiza\c c\~ao e Justificativa}

A aplica\c c\~ao de t\'ecnicas de aprendizado de m\'aquina ao problema de diagn\'ostico e predi\c c\~ao de doen\c cas cardiovasculares tem atra\'ido aten\c c\~ao crescente tanto na literatura cient\'ifica internacional quanto no contexto acad\^emico brasileiro. Estudos recentes demonstram que algoritmos supervisionados, incluindo \'Arvores de Decis\~ao, M\'aquinas de Vetores de Suporte, Florestas Aleat\'orias, K-Vizinhos Mais Pr\'oximos e diversas arquiteturas de redes neurais, s\~ao capazes de atingir n\'iveis de acur\'acia competitivos com m\'etodos diagn\'osticos tradicionais quando aplicados a conjuntos de dados cl\'inicos estruturados \cite{mesquita2020,ravani2024}. Contudo, a rela\c c\~ao entre complexidade do modelo e desempenho preditivo n\~ao \'e linear, e a quest\~ao de em que circunst\^ancias a complexidade adicional de um modelo de aprendizado de m\'aquina se traduz em ganho cl\'inico real permanece em aberto.

O presente trabalho situa-se precisamente nessa quest\~ao, investigando-a no contexto espec\'ifico das Redes Neurais Artificiais (RNAs). As RNAs constituem uma das fam\'ilias de modelos de maior expans\~ao em aplica\c c\~oes biom\'edicas, e dentro delas existem arquiteturas de complexidades muito distintas. No extremo mais simples, o ADALINE Log\'istico (\textit{Adaptive Linear Neuron}, ADALINE, com sa\'ida log\'istica) \'e um modelo de neur\^onio \'unico que aplica a fun\c c\~ao sigmoide em cada passo de treinamento para obter a sa\'ida preditiva e atualiza os pesos por uma regra proporcional ao erro de sa\'ida (regra delta/LMS), sem propagar o gradiente pela derivada da sigmoide, o que o distingue da retropropaga\c c\~ao plena do MLP, sendo capaz de separar classes linearmente separ\'aveis com base em combina\c c\~oes lineares dos atributos de entrada \cite{braga2007}. No extremo de maior capacidade representacional, o Perceptron Multicamadas (MLP) incorpora uma ou mais camadas ocultas de neur\^onios com ativa\c c\~oes n\~ao lineares, o que lhe confere, segundo o Teorema da Aproxima\c c\~ao Universal \cite{cybenko1989,hornik1991}, a capacidade te\'orica de aprender qualquer fun\c c\~ao cont\'inua em dom\'inios compactos \cite{braga2007,silva2016}.

A pergunta central deste trabalho \'e: \textit{para dados cl\'inicos tabulares de cardiopatas, o MLP produz resultados expressivamente superiores ao ADALINE?} A resposta a essa pergunta tem implica\c c\~oes pr\'aticas relevantes: se o ADALINE for suficiente, seu uso \'e prefer\'ivel pela menor sensibilidade a hiperpar\^ametros, maior interpretabilidade e estabilidade de treinamento; se o MLP for significativamente superior, sua complexidade adicional se justifica. Al\'em disso, a compara\c c\~ao expl\'icita entre as vers\~oes implementadas do zero (NumPy) e as vers\~oes com biblioteca (scikit-learn) permite verificar a corretude das implementa\c c\~oes pr\'oprias e quantificar a diferen\c ca de desempenho introduzida pelas otimiza\c c\~oes internas das bibliotecas.

A escolha de dois conjuntos de dados com perfis cl\'inicos distintos refor\c ca a robustez da investiga\c c\~ao. O Banco 1 \'e o UCI Heart Disease Dataset \cite{fedesoriano2021}, com 918 pacientes e vari\'aveis de triagem diagn\'ostica (press\~ao arterial, colesterol, ECG, angina induzida por esfor\c co, entre outras). O Banco 2 \'e o Heart Failure Clinical Records Dataset \cite{chicco2020}, com 299 pacientes j\'a diagnosticados com insufici\^encia card\'iaca, cujo objetivo \'e a predi\c c\~ao de mortalidade com base em biomarcadores de progn\'ostico (fra\c c\~ao de eje\c c\~ao, creatinina s\'erica, s\'odio s\'erico, entre outros). Avaliar os modelos em dois contextos cl\'inicos diferentes, triagem diagn\'ostica e progn\'ostico de mortalidade, permite identificar se os padr\~oes de desempenho observados s\~ao espec\'ificos de um conjunto de dados ou se refletem propriedades gerais das arquiteturas sob investiga\c c\~ao. Cabe destacar que o pr\'e-processamento adotado converte todas as vari\'aveis, cont\'inuas e categ\'oricas, em representa\c c\~oes bin\'arias com base nos limiares das diretrizes cl\'inicas brasileiras, op\c c\~ao adotada com o intuito de alinhar a representa\c c\~ao \`a categoriza\c c\~ao m\'edica oficial e favorecer a interpretabilidade dos modelos, cujo efeito emp\'irico sobre o desempenho \'e investigado no experimento complementar da Se\c c\~ao~\ref{sec:resultados_bruto}.

\section{Objetivos}

\subsection{Objetivo Geral}

Comparar o desempenho preditivo do ADALINE Log\'istico e do Perceptron Multicamadas na classifica\c c\~ao bin\'aria de doen\c cas card\'iacas, avaliando tanto a acur\'acia global quanto o n\'umero de Falsos Negativos como crit\'erio cl\'inico de prioridade.

\subsection{Objetivos Espec\'ificos}

\begin{itemize}
  \item Implementar o ADALINE Log\'istico e o Perceptron Multicamadas do zero em Python com NumPy, formalizando matematicamente as regras de atualiza\c c\~ao de pesos de cada arquitetura.
  \item Validar as implementa\c c\~oes pr\'oprias por correspond\^encia com as vers\~oes equivalentes da biblioteca scikit-learn.
  \item Pr\'e-processar os dois conjuntos de dados com base nos limiares das diretrizes cl\'inicas brasileiras de cardiologia, convertendo vari\'aveis cont\'inuas e categ\'oricas em representa\c c\~oes bin\'arias interpret\'aveis.
  \item Avaliar os quatro modelos sob tr\^es configura\c c\~oes sistem\'aticas de treinamento, totalizando 24 experimentos independentes com semente aleat\'oria fixada.
  \item Analisar os resultados sob dois crit\'erios complementares: acur\'acia global e n\'umero de Falsos Negativos, justificando a prioridade cl\'inica do segundo crit\'erio no contexto de diagn\'ostico card\'iaco.
  \item Identificar os fatores estruturais respons\'aveis pelas diferen\c cas de desempenho observadas entre arquiteturas e configura\c c\~oes de treinamento, incluindo o efeito da binariza\c c\~ao das vari\'aveis avaliado em experimento complementar com dados brutos.
\end{itemize}

\section{Organiza\c c\~ao do Trabalho}

O trabalho est\'a organizado em seis cap\'itulos al\'em desta Introdu\c c\~ao: o Cap\'itulo~\ref{cap:fundamentacao} apresenta os fundamentos te\'oricos do ADALINE Log\'istico e do MLP; o Cap\'itulo~\ref{cap:trabalhos} revisa tr\^es trabalhos relacionados em aplica\c c\~oes cardiovasculares; o Cap\'itulo~\ref{cap:metodologia} descreve os conjuntos de dados, o pr\'e-processamento e as implementa\c c\~oes; o Cap\'itulo~\ref{cap:resultados} apresenta os resultados dos 24 experimentos e o experimento complementar com dados brutos; o Cap\'itulo~\ref{cap:discussao} interpreta os resultados \`a luz da teoria; e o Cap\'itulo~\ref{cap:conclusao} sintetiza as principais constata\c c\~oes e aponta perspectivas de continuidade.
